% Corrected French for printed p.220.
% Three exact print-level repairs are recorded in SOURCE_CORRECTIONS.jsonl.
\clearpage
\SourcePage{220}

\LayerTarget{tohoku:transition:section-5.7-functorial-formulas}
\LayerTarget{tohoku:equation:5.7.3}
\begin{equation}
\Gamma_\Phi^G=\Gamma_\Psi f_*^G=\Gamma^G\Gamma_\Phi.
\tag{5.7.3}
\end{equation}

\LayerTarget{tohoku:spectral-sequence:5.7-i-supported-equivariant}
Comme $f_*^G$ transforme objets injectifs de $C^{X(G)}$ en faisceaux
flasques, donc $\Gamma_\Psi$-acycliques, la première formule donne encore
\textit{un foncteur spectral cohomologique sur $C^{X(G)}$, aboutissant au
foncteur gradué $(H_\Phi^m(X;G,A))$, et de terme initial}

\LayerTarget{tohoku:equation:5.7.4}
\begin{equation}
\mathrm I_2^{p,q}(A)=H_\Psi^p\bigl(Y,\mathbf H^q(G,A)\bigr).
\tag{5.7.4}
\end{equation}

\LayerTarget{tohoku:obstruction:second-spectral-sequence-with-supports}
La seconde suite spectrale du théorème 5.2.1 ne se généralise plus telle
quelle, car on ne pourra pas dire en général que si $B$ est un objet injectif
de $C^{X(G)}$, $\Gamma_\Phi(B)$ soit un $G$-module $\Gamma^G$-acyclique. Mais
\textit{supposons que tout ensemble de $\Psi$ ait un voisinage qui soit dans
$\Psi$.}

\LayerTarget{tohoku:construction:extension-by-zero-support-limit}
Pour tout ouvert $U\subset Y$, soit $A_U=A_{f^{-1}(U)}$ le $G$-faisceau
abélien sur $X$ qui est nul dans $C f^{-1}(U)$ et coïncide avec $A$ sur
$f^{-1}(U)$. On a
$\Gamma_\Phi(A)=\varinjlim\Gamma_X(A_U)$, où la limite inductive est prise
suivant l'ensemble ordonné filtrant des parties ouvertes $U$ de $Y$ telles
que $\overline U\in\Psi$. Comme les applications
$\Gamma_X(A_U)\to\Gamma_X(A_V)$ pour $U\subset V$ sont injectives, le foncteur
$\Gamma^G$ passe à la limite inductive, et on obtient une identité
fonctorielle :

\LayerTarget{tohoku:equation:5.7.5}
\begin{equation}
\Gamma_\Phi^G(A)=\varinjlim\Gamma_X^G(A_U)
\tag{5.7.5}
\end{equation}
d'où, comme pour le corollaire 1 à prop. 3.10.1 :

\LayerTarget{tohoku:equation:5.7.6}
\begin{equation}
H_\Phi^m(X;G,A)=\varinjlim H^m(X;G,A_U).
\tag{5.7.6}
\end{equation}

\LayerTarget{tohoku:spectral-sequence:ii-a-u}
Or pour tout $U$, $H^*(X;G,A_U)$ est l'aboutissement d'une suite spectrale
cohomologique $\mathrm {II}(A,U)$ de terme initial

\LayerTarget{tohoku:equation:ii-a-u-initial-term}
\[
\mathrm {II}(A,U)_2^{p,q}=H^p\bigl(G,H^q(X,A_U)\bigr).
\]

\LayerTarget{tohoku:construction:inductive-system-of-spectral-sequences}
et pour $U$ variable, ces suites spectrales forment un système inductif. On
conclut alors de (5.7.6) que $H_\Phi^*(X;G,A)$ est \textit{l'aboutissement
d'une suite spectrale cohomologique de terme initial}

\LayerTarget{tohoku:equation:5.7.7}
\begin{equation}
\mathrm {II}_2^{p,q}(A)=\varinjlim_U H^p\bigl(G,H^q(X,A_U)\bigr).
\tag{5.7.7}
\end{equation}

\LayerTarget{tohoku:comparison:cohomology-and-filtered-limit}
Bien entendu, cette suite spectrale est même un foncteur spectral en $A$.
Dans les cas les plus importants (que nous allons expliciter ci-dessous) on
pourra échanger les signes $H^p(G,-)$ et $\varinjlim$, et alors on obtient,
puisque $\varinjlim H^q(X,A_U)=H_\Phi^q(X,A)$ (corollaire 1 de la proposition
3.10.1), la forme usuelle du terme initial :

\LayerTarget{tohoku:equation:5.7.7-bis}
\begin{equation}
\mathrm {II}_2^{p,q}=H^p\bigl(G,H_\Phi^q(X,A)\bigr)
\tag{5.7.7 bis}
\end{equation}

\LayerTarget{tohoku:conditions:validity-5.7.7-bis}
\noindent\textit{Conditions de validité de la formule (5.7.7 bis).}

\LayerTarget{tohoku:condition:5.7.7-bis-finite-group}
\noindent a) $G$ est un groupe fini.

\LayerTarget{tohoku:condition:5.7.7-bis-trivial-action}
\noindent b) $G$ opère trivialement dans les $H^q(X,A_U)$, et les
$H_p(G,\mathbf Z)$ sont de type fini.

\LayerTarget{tohoku:condition:5.7.7-bis-essentially-constant}
\noindent c) Le système inductif $(H^q(X,A_U))$ est
\og essentiellement constant\fg{} pour tout $q$, i.~e. on peut trouver un
système cofinal $(U_i)$ tel que les homomorphismes
$H^q(X,A_{U_i})\to H^q(X,A_{U_j})$ soient des isomorphismes.

\noindent\textit{Qualification éditoriale.} L'hypothèse de voisinage ci-dessus
ne démontre que (5.7.5). L'identité de limite inductive dérivée (5.7.6) et
l'aboutissement affirmé avant (5.7.7) exigent en outre que, pour tout $U$
admissible, le foncteur exact de prolongement par zéro
$E_U:A\longmapsto A_U$ transforme les $G$-faisceaux abéliens injectifs en
objets $\Gamma_X^G$-acycliques. La comparaison avec la cohomologie ordinaire
utilisée dans (5.7.7 bis) exige aussi que $\Phi$ soit paracompactifiante. À
défaut de ces hypothèses supplémentaires, aucune identité de limite inductive
dérivée ni aucun aboutissement n'est affirmé ici.

\TohokuContinuation{221}
